\documentclass[10pt]{article}
\usepackage[margin=0.9in]{geometry}
\usepackage[T1]{fontenc}
\usepackage{lmodern}
\usepackage{amsmath,amssymb}
\usepackage{booktabs}
\usepackage{tabularx}
\usepackage{hyperref}
\usepackage{enumitem}
\usepackage{float}
\setlist{nosep}
\setlength{\parskip}{0.4em}
\setlength{\parindent}{0pt}
\newcommand{\BC}{{\mdseries\textsc{BoundClaw}}}
\newcommand{\M}[1]{\textit{[Methodology #1]}}
\newcolumntype{L}[1]{>{\raggedright\arraybackslash}p{#1}}
\newcolumntype{Y}{>{\raggedright\arraybackslash}X}

\title{Policy Inference in \BC: the Machine-Learning View}
\date{Companion to the evaluation methodology, \today}

\begin{document}
\maketitle

\noindent This note restates the evaluation methodology (\texttt{methodology.pdf}) as a machine-learning
problem. It changes nothing in that plan; it only translates it. Each section names the part of the
methodology it covers, in the form \M{\S1}.

\section{The learning problem}
\M{\S1, Table 1}

\paragraph{Setting.} One agent hands a subtask to another agent. Before the second agent starts, a model
decides which of the first agent's permissions to pass along. Passing too many is a security cost; passing
too few can make the task fail.

\paragraph{Input.} Text, plus a list of candidates:
\begin{itemize}
  \item the \emph{root task}: the user's original request;
  \item the \emph{subtask}: the text the first agent wrote for the second;
  \item context: who receives the subtask, earlier conversation, and an optional ``manifest'' in which the
    receiver's operator lists what it usually needs (a hint that may be wrong);
  \item a numbered \emph{candidate list}, built by a fixed program, not a model. Each candidate is one
    permission the first agent already holds, such as ``write objects under \texttt{staging/build/}'',
    ``connect to \texttt{docs.example}'' or ``mount the \texttt{workspace} folder read-only''. When the text names a
    specific file inside a broader permission, the narrower permission is added as a further candidate.
\end{itemize}

\paragraph{Output.} One label per candidate, in three classes: \emph{required} (granted now),
\emph{plausible} (may be granted later, see below) and \emph{unneeded}. The model outputs only candidate
indices in fixed JSON. It is a \textbf{set-selection task over a closed list}, which is per-item three-class
classification conditioned on the text and on the whole list.

\paragraph{What the model cannot do.} It cannot name a permission outside the list, so its worst possible
output (select everything) equals passing every permission along, which is the no-model baseline. Rules
outside the model enforce this. The model therefore affects \emph{how much} permission is passed, never
whether the system's hard limits hold. We evaluate it as a component whose errors are bounded.

\paragraph{The ``plausible'' class is a second chance.} If the receiver later hits a permission it lacks, the
system may grant it, one permission per request and at most $k=3$ requests per subtask (refused requests
count). It grants one only if the model put that candidate in \emph{plausible}. So \emph{required} is the
point prediction and \emph{plausible} is a backup set. A large backup set fixes missed permissions but lets an
attacker who triggers a request get more; a small one is safe but fixes less. Folder (volume) permissions get
no second chance, so missing one fails the task.

\paragraph{Two passes, keep the minimum.} Each candidate is scored twice: once against the root task and once
against the subtask. The lower score is kept. This is a conservative two-view ensemble: if an attacker edits
the subtask to ask for extra permission, the root task, which the attacker did not write, still scores it low.

\paragraph{Abstention.} A timeout, an error, or output that breaks the format (an unknown key, an index out of
range, trailing text) refuses the whole handoff. There is no default answer; invalid output is never
repaired.

\paragraph{Receiver limit.} Each receiver has a maximum it will accept (its ``ceiling'') that the model never
sees. If the selected network or folder permissions exceed it, the handoff is rejected. Over-selection can
therefore cause failures too, not only security cost.

\begin{table}[H]\centering\small
\begin{tabularx}{\linewidth}{L{0.3\linewidth}Y}
\toprule
Security term (methodology) & Meaning here \\
\midrule
Helper & The model \\
Delegator, receiver & The agent handing off the subtask, and the agent receiving it \\
$C$ (capability) & Permissions on services, such as reading or writing objects under a storage path \\
$N$, \emph{plain} & Network hosts the receiver may contact; \emph{plain} hosts are reached without credentials \\
$D$ (volumes) & Folders mounted into the receiver, read-only or read-write \\
Proposal & The \emph{required} class \\
Admission set $A$ & The \emph{plausible} class \\
Recovery & The second-chance request at run time \\
G1, checker & Hard rules outside the model; they bound every output \\
Ceiling $P_{\mathrm{ceil}}$ & The receiver's hidden maximum; exceeding it rejects the handoff \\
Fail closed, refusal & Abstention: no output means no handoff \\
Injection & Prompt injection: text planted in the input to make the model select extra permissions \\
\bottomrule
\end{tabularx}
\caption{Terms used in the methodology and their meaning in this note.}
\end{table}

\section{Why this design: what the pilot study showed}
\M{\S2, Table 2}

We first ran a broad pilot: 23 model configurations of four kinds on an earlier version of the task. The pilot
numbers are \emph{not} reported as results; they only explain the design. Main lessons:
\begin{itemize}
  \item \textbf{The backup set is the key safety lever.} Without the plausible gate, every planted attack could
    get its permission through the second chance. With it, 13\% could, and 97\% of tasks still finished.
  \item \textbf{Models must output three classes.} Models whose scores were nearly binary could not separate
    ``required'' from ``plausible'', so their backup set gave no protection.
  \item \textbf{Mid-size listwise LLMs were the most precise.} Qwen3-4B, 8B and 30B-A3B passed 6--9\% of
    unneeded permissions; the largest model (gpt-oss-120b) passed 14\%.
  \item \textbf{Fast pointwise models need training.} A 0.6B reranker answers in 28\,ms but, untrained, passed 40\%
    of unneeded permissions, against 51\% for a fair coin. They lack task knowledge, which training can supply.
  \item \textbf{The manifest hint helps} (23\% to 9\% unneeded for Qwen3-8B), so a poisoned hint is an attack worth testing.
  \item \textbf{Top-10 retrieval does not fit.} CAPMAS, the closest prior system, keeps a fixed top 10 and passed 74\%.
  \item \textbf{Local open-weight models only}, because the model reads possibly confidential text and open
    weights allow white-box attacks.
\end{itemize}

\section{Research questions}
\M{\S3}
\begin{description}
  \item[RQ1 Precision and recall.] How many unneeded permissions does the model remove, how many needed ones does
    it miss, and how often does over-selection hit the receiver limit?
  \item[RQ2 The backup set as a gate.] Does restricting second chances to the plausible class stop attacks,
    compared with allowing any permission the first agent holds, and what does it cost honest tasks?
  \item[RQ3 Adversarial robustness.] How much of the RQ1 gain survives (a) planted prompt injections, (b)
    injection styles never seen in training (distribution shift), (c) a white-box adaptive attacker, (d) an
    edited subtask, with and without the two-pass minimum, and (e) a poisoned manifest hint?
  \item[RQ4 Inference cost.] Latency per handoff (two passes) and abstention rate.
  \item[RQ5 Fine-tuning.] Can small, fast models match the best prompted model after training, including on task
    families and attack styles never seen in training (out-of-distribution)?
\end{description}

\section{Models}
\M{\S4, Table 3}
\begin{table}[H]\centering\small
\begin{tabularx}{\linewidth}{L{0.24\linewidth}L{0.2\linewidth}L{0.2\linewidth}Y}
\toprule
Model & Architecture & Variants & Why \\
\midrule
Qwen3-8B & listwise LLM & prompted, fine-tuned & Best pilot trade-off; the reference model \\
Qwen3-4B, 1.7B & listwise LLM & prompted, fine-tuned & Smaller sizes \\
Qwen3-30B-A3B & mixture-of-experts LLM & prompted, QLoRA & Fastest precise LLM; must learn the plausible class \\
Qwen3-Reranker-0.6B, 8B & pointwise cross-encoder & prompted, fine-tuned & Fastest model, and best pointwise one \\
CAPMAS & bi-encoder, contrastive & as published; retrained & Closest prior system \\
gpt-oss-120b & listwise LLM & prompted & Upper reference, too slow to deploy \\
Baselines & -- & -- & Select all; manifest only; keyword match; coin flips at 25/50/75\% \\
\bottomrule
\end{tabularx}
\end{table}
Listwise LLMs read the whole candidate list and return index lists. Pointwise models score each candidate
alone, and two thresholds turn scores into the three classes (tuned on the development split).

\section{Data}
\M{\S5}
\paragraph{Evaluation set (human-labeled).} 408 tasks from 17 templates: 288 test and 120 development, split
\emph{by template} so no template leaks across. Each template varies over 4 name sets, 3 levels of how many
extra permissions the first agent holds (how many distractor candidates the model sees), and with or without
one prompt injection. Injections appear in the subtask or context, never in the root task. Two annotators label
every candidate independently and resolve disagreements; we report agreement (Cohen's $\kappa$). The test set is
frozen by checksum before any model runs.

\paragraph{Synthetic training set.} No public dataset fits, so we generate one.
\begin{itemize}
  \item \emph{Size and split.} 2,000 tasks in 40 task families: 1,500 train and 300 validation from 32 families,
    and 200 test tasks from 8 families held out entirely (out-of-distribution test). Size follows the closest
    prior work. Tasks sharing any 8-word sequence or name with the evaluation set are removed (decontamination).
  \item \emph{Generation.} A local LLM writes the root task, the subtask, and the API calls that complete the
    subtask, optionally with a second valid plan.
  \item \emph{Programmatic labels.} The calls run against mock services that log which permissions they used.
    The \emph{required} label is the smallest candidate set covering that log; \emph{plausible} is what the second
    plan adds. No model writes labels. Checked: this labeler reproduces the human labels on all 408 evaluation tasks.
  \item \emph{Filtering.} Format checks, mock execution, and two judge LLMs (not the writer) that must both
    accept the task; near-duplicates are removed. We report how many tasks survive each step.
  \item \emph{Distractors and injections.} As in the evaluation set. Half the tasks carry an injection with an
    unchanged label, so the model learns to ignore it.
  \item \emph{Human audit.} Two people label 93 random tasks without seeing the generated labels. The set passes
    if at most one generated label is wrong, which bounds the label error rate below 5\% with 95\% confidence.
    Otherwise we fix the generator and redo it.
\end{itemize}

\paragraph{External benchmarks.} FORTIS and the mostargate prompts, from other groups, for transfer. Test only.

\paragraph{System test.} The models run inside the real \BC{} prototype on object-storage workflows with scripted
and live agents, and an automatic judge that reads system logs, not the agent's reply.

\paragraph{Attack suite.} The five RQ3 attacks. The adaptive attacker knows the format and has the weights, and an
attacker LLM searches for injected text that gets a target permission selected, under a fixed query budget; we
report success against budget.

\section{Training}
\M{\S6}
Training cannot make the system safer in a hard sense; the external rules do that. It can make small models
more precise and faster.
\begin{itemize}
  \item \emph{Targets.} Pointwise models: needed or not per candidate. Listwise LLMs: the three-class index output.
    CAPMAS: its own contrastive objective.
  \item \emph{Method.} Full fine-tuning (all weights) up to 8B, so trained and untrained versions differ only by
    training. QLoRA for the 30B model, compared only at the same precision.
  \item \emph{Settings.} Cosine schedule, 3\% warm-up, 2 epochs, loss on output tokens only; learning rate from
    $\{5\!\times\!10^{-6}, 10^{-5}, 2\!\times\!10^{-5}\}$ chosen on validation; 3 seeds.
  \item \emph{No leakage.} Only synthetic training and validation data are used for training and model selection.
    The held-out families are scored once.
\end{itemize}

\section{Metrics and statistics}
\M{\S7}
Permissions that cover many items (a whole path prefix) count for every item they cover.
\begin{itemize}
  \item \textbf{Excess}: share of the unneeded permissions that the model still grants (selecting everything
    scores 1). A false-positive rate normalized by the baseline.
  \item \textbf{Missing}: share of needed permissions not granted ($1-$recall). \textbf{Exact match}: everything right.
  \item \textbf{Limit rejections}: share of handoffs rejected by the receiver's hidden maximum.
  \item \textbf{Attack success after second chances}: share of attacks whose target is granted at once or within 3
    requests (also 1 and 5), with and without the plausible gate.
  \item \textbf{Task completion}, counting abstentions, rejections and missing folders as failures.
  \item \textbf{Gain retained under attack}: 1 means the attack gained nothing, 0 means the model was pushed back to
    selecting everything.
  \item \textbf{Latency}: median and 95th percentile, one model per GPU.
\end{itemize}
Comparisons and pass rules are fixed before results are seen. Confidence intervals use bootstrap resampling
by template or family, since tasks from one template are correlated. Paired yes/no outcomes use McNemar's
test; multiple comparisons use Holm correction; trained models report mean and spread over seeds.

\section{Open questions}
\M{\S8}
Decisions the systems team still has to make: which services besides object storage stay in the task; which
receivers are local or remote and what each one's limit is; the request cap (we propose 3) and the model
timeout; which hosts are reached without credentials; and whether system tests may log model outputs.

\end{document}
